% ================================================================
% TAILORING
% ================================================================

\section{Tailoring}
\label{sec:tailoring}

Lo standard ISO/IEC~12207:1995\textsubscript{G} prevede la possibilità di adattare
(\textit{tailoring}\textsubscript{G}) i processi del ciclo di vita software al contesto
specifico del progetto.

Il gruppo \textit{Synergo Unit} adotta un approccio incrementale orientato al
miglioramento continuo, introducendo progressivamente processi e pratiche
in funzione dell'evoluzione del progetto e delle baseline\textsubscript{G} previste dal
regolamento didattico.

Il tailoring adottato tiene conto dei seguenti fattori contestuali:

\begin{itemize}
    \item \textbf{Baseline attuale}: RTB\textsubscript{G}; il gruppo si trova nella fase di consolidamento dei requisiti e della documentazione tecnica e gestionale richiesta per la Requirements \& Technology Baseline\textsubscript{G};

    \item \textbf{Dimensione del gruppo}: sei membri con rotazione periodica e documentata dei ruoli;

    \item \textbf{Natura didattica del progetto}: il progetto viene svolto in ambito accademico, con vincoli temporali, organizzativi e tecnologici definiti dal regolamento del corso;

    \item \textbf{Approccio incrementale}: il gruppo introduce esclusivamente processi effettivamente utilizzati e verificabili.
\end{itemize}

Il presente documento è soggetto ad aggiornamento progressivo secondo il principio \textit{just-in-time}\textsubscript{G}: ogni nuova attività o pratica introdotta nel progetto deve essere normata prima della sua applicazione operativa.

% ------------------------------------------------

\subsection{Processi attivati nella fase CC}

La seguente tabella riporta i processi attivati durante la fase di candidatura\textsubscript{G} (CC), mantenuti operativi anche nelle fasi successive del progetto.

\bigskip

\begin{tabularx}{\textwidth}{l l X}
    \toprule
    \textbf{Categoria} & \textbf{Processo} & \textbf{Sezione} \\
    \midrule

    Primario
    & Fornitura
    & \ref{sec:fornitura} \\

    Di supporto
    & Gestione della documentazione
    & \ref{sec:documentazione} \\

    Di supporto
    & Gestione della configurazione
    & \ref{sec:configurazione} \\

    Di supporto
    & Verifica documentale
    & \ref{sec:verifica} \\

    Organizzativo
    & Pianificazione e gestione delle attività
    & \ref{sec:gestione} \\

    Organizzativo
    & Gestione dell'infrastruttura operativa
    & \ref{sec:infrastruttura} \\

    \bottomrule
\end{tabularx}

% ------------------------------------------------

\subsection{Processi attivati nella fase RTB}

Con l'avanzamento alla baseline RTB vengono introdotti ulteriori processi necessari al consolidamento dei requisiti, della qualità e delle attività di progettazione.

I processi precedentemente attivati rimangono operativi.

\bigskip

\begin{tabularx}{\textwidth}{l l X}
    \toprule
    \textbf{Categoria} & \textbf{Processo} & \textbf{Sezione} \\
    \midrule

    Primario
    & Analisi dei requisiti
    & \ref{sec:sviluppo} \\

    Primario
    & Analisi dello stato dell'arte
    & \ref{sec:sviluppo} \\

    Primario
    & Progettazione architetturale preliminare
    & \ref{sec:sviluppo} \\

    Primario
    & Proof of Concept
    & \ref{sec:sviluppo} \\

    Di supporto
    & Garanzia della qualità
    & \ref{sec:qualita} \\

    Di supporto
    & Gestione dei rischi
    & \ref{sec:rischi} \\

    Di supporto
    & Utilizzo di strumenti AI
    & \ref{sec:ai} \\

    \bottomrule
\end{tabularx}

% ------------------------------------------------

\subsection{Processi pianificati per la fase PB}

Alcuni processi risultano attualmente non necessari oppure ancora in consolidamento.

Tali processi verranno rivalutati e, ove necessario, introdotti durante la Product Baseline\textsubscript{G} (PB).

Tra i principali processi pianificati rientrano:

\begin{itemize}
    \item integrazione continua;
    \item testing automatico;
    \item deployment automatizzato;
    \item gestione avanzata delle release software;
    \item metriche avanzate di prodotto;
    \item Architecture Decision Records (ADR)\textsubscript{G}.
\end{itemize}

L'introduzione di nuovi processi comporterà l'aggiornamento incrementale del presente documento.